% =====================================================================
\subsection{Workflow di provisioning multi-dominio}
\label{sec:workflow}
% =====================================================================

La clausola 6 lega le informazioni al ciclo di vita del servizio. Quando la
rete QKD viene introdotta, l'orchestrator ne acquisisce la topologia
(§5.1); durante l'esercizio riceve, su richiesta o periodicamente, stato e
inventario (§5.2); dopo l'attivazione di un servizio ne monitora lo stato
(§5.3); in ogni momento riceve notifiche (§5.5).

La §6.2 descrive il caso completo (Fig.~\ref{fig:seq018}): un orchestrator
coordina il controller QKD e il controller OTN, che non comunicano tra loro
direttamente.

% Figura: provisioning end-to-end QKD + OTN (adattata da GS QKD 018, Fig. 2)
\begin{figure}[t]
\centering
\begin{tikzpicture}[font=\scriptsize]
  \def\xA{0}\def\xB{4.0}\def\xC{8.4}\def\xD{12.2}
  \node[actor] at (\xA,0) {Amministratore};
  \node[actor] at (\xB,0) {Orchestrator};
  \node[actor,fill=green!8] at (\xC,0) {Controller QKD};
  \node[actor,fill=orange!10] at (\xD,0) {Controller OTN};
  \foreach \x in {\xA,\xB,\xC,\xD} \draw[gray] (\x,-0.4) -- (\x,-9.0);
  \draw[msg] (\xA,-1.0) -- node[above]{richiesta servizio (SAE, QoS)} (\xB,-1.0);
  \draw[msg] (\xB,-1.5) -- ++(0.4,0) -- ++(0,-0.4) -- (\xB,-1.9);
  \node[right,align=left] at (\xB+0.45,-1.7) {address matching\\nodo OTN $\leftrightarrow$ nodo QKD};
  \draw[blue!60,rounded corners] (\xB-0.3,-2.55) rectangle (\xC+0.3,-3.85);
  \node[blue!60,anchor=south west,font=\scriptsize\itshape] at (\xB-0.3,-2.55) {ciclo ripetibile};
  \draw[msg] (\xB,-3.0) -- node[above]{creazione service link + vincoli} (\xC,-3.0);
  \draw[ret] (\xC,-3.6) -- node[above]{\texttt{svc\_link\_id}, cammini candidati} (\xB,-3.6);
  \draw[msg] (\xB,-4.6) -- node[above]{deploy del cammino scelto (\texttt{path\_id})} (\xC,-4.6);
  \draw[msg] (\xC,-5.0) -- ++(0.4,0) -- ++(0,-0.4) -- (\xC,-5.4);
  \node[right,align=left] at (\xC+0.45,-5.2) {key relay lungo\\il cammino};
  \draw[ret] (\xC,-6.1) -- node[above]{chiavi pronte} (\xB,-6.1);
  \draw[msg] (\xB,-6.9) -- node[above]{provisioning applicazione (\texttt{qkd\_app})} (\xC,-6.9);
  \draw[ret] (\xC,-7.6) -- node[above]{\texttt{sdqkdn\_application\_deployed}} (\xB,-7.6);
  \draw[msg] (\xB,-8.6) -- node[above,fill=white,inner sep=1pt]{i nodi OTN possono prelevare le chiavi dai nodi QKD} (\xD,-8.6);
\end{tikzpicture}
\caption[Provisioning end-to-end di un servizio QKD tra rete QKD e OTN]{Provisioning end-to-end di un servizio QKD tra rete QKD e OTN
(rielaborazione semplificata di~\cite[Fig.~2]{qkd018}). L'interazione con
il controller OTN è fuori dallo scope della specifica.}
\label{fig:seq018}
\end{figure}


I passi sono:
\begin{enumerate}
  \item l'amministratore chiede all'orchestrator un servizio, indicando i
        SAE nei nodi OTN e la loro QoS;
  \item l'orchestrator individua i controller coinvolti ed effettua
        l'address matching; conosce già gli indirizzi dei nodi QKD grazie
        a topologia e inventario;
  \item chiede al controller QKD la creazione di un service link e il
        calcolo dei cammini, con i vincoli desiderati;
  \item il controller assegna un \texttt{svc\_link\_id} e restituisce i
        cammini candidati, con una priorità per ciascuno; il passo può
        ripetersi con vincoli diversi;
  \item l'orchestrator sceglie un cammino e ne chiede il deploy; i nodi
        generano e inoltrano chiavi lungo il cammino finché i due estremi
        sono pronti;
  \item ricevuta la conferma, l'orchestrator richiede il provisioning
        dell'applicazione sul service link e avvisa il controller OTN che i
        nodi OTN possono iniziare a ricevere chiavi dai nodi QKD indicati.
\end{enumerate}
Si noti il parallelo con il workflow del link virtuale della 015: lì il
controller calcolava e configurava il cammino da solo, qui lo
\emph{propone} e attende la scelta dell'orchestrator. Con più domini QKD
l'orchestrator ripete i passi 3–5 per ciascun dominio lungo il cammino
inter-dominio scelto, collegando i segmenti attraverso i nodi di bordo: è
ciò che fanno Sim et al.\ (§\ref{sec:sim}).

% =====================================================================
\subsection{Protocolli e sicurezza}
\label{sec:protocolli}
% =====================================================================

\paragraph{Trasporto (§8).}
Come la 015, lo standard separa il modello dei dati (YANG~1.1) dal
protocollo di trasporto, che può essere NETCONF~\cite{rfc6241} o
RESTCONF~\cite{rfc8040}. Con RESTCONF ogni risorsa è un URI, la lettura usa
\texttt{GET} e la configurazione \texttt{POST}, \texttt{PUT},
\texttt{PATCH} e \texttt{DELETE}, con codifica XML o JSON; le notifiche si
ricevono sottoscrivendo i relativi stream. Lo standard osserva che il
paradigma REST, senza stato, si presta bene a questo tipo di API. A titolo
di esempio, la lettura della topologia e l'invocazione di una RPC avrebbero
la forma:
\begin{lstlisting}
GET  /restconf/data/etsi-qkd-sdn-topology:qkd_network
POST /restconf/operations/etsi-qkd-sdn-connectivity:create-service-link-id-for-app
\end{lstlisting}

\paragraph{Sicurezza (§7).}
Il controller ha una vista centrale della rete e ne passa un sottoinsieme
all'orchestrator; le misure di protezione dell'interfaccia (autenticazione,
cifratura del canale, per esempio TLS sotto RESTCONF) sono lasciate
all'implementazione. Il principio fondamentale, ereditato dalla GS~QKD~015,
è che \textbf{le chiavi derivate da QKD non transitano mai né per il
controller né per l'orchestrator}: nei moduli YANG compaiono solo tassi di
chiave e l'indirizzo del KME. Un orchestrator compromesso può quindi
causare disservizi (cammini sbagliati, risorse esaurite, denial of service)
ma non leggere le chiavi. Lo standard raccomanda procedure accurate di
progettazione, installazione e personalizzazione per mitigare richieste
malevole e attacchi ai canali di controllo.
